\documentclass[../main]{subfiles}
\begin{document}
\chapter{Learning Theory}
\begin{introduction}
    \item In-Sample Error
    \item Out-of-Sample Error
    \item Generalization Error
    \item PAC
    \item Shelah's Lemma
    \item VC-Dim
\end{introduction}
\section{Learning Theory}
Learning theory forms the foundation of machine learning. It explains why a model can generalize to data that did not appear during training. We may draw an analogy here:
\begin{itemize}
\item Training a model $\to$ Taking a course;
\item Training process $\to$ Doing exercises or homework;
\item Testing phase $\to$ Taking an exam.
\end{itemize}
From our own learning experience, we know that students who perform well during practice do not always achieve high scores in exams. Likewise, a model that fits the training data well may still perform poorly on unseen data. This motivates us to introduce theoretical concepts that can quantitatively describe and explain such phenomena.

\begin{definition}[In-sample Error]
    The \emph{in-sample error} $E_{\rm in}$ is defined as the empirical error rate of a hypothesis evaluated on the training dataset.
    \end{definition}
    
    Let $h \in \mathcal{H}$ denote a hypothesis or model, for example a linear classifier $h(x) := \mathrm{sign}(w^\top x + b)$.  
    Then the in-sample error can be formally written as
    \begin{equation}
        E_{\rm in} = \frac{1}{n}\sum_{i \in [n]} \mathbb{1}\!\big(h(x_i) \neq y_i\big),
        \qquad D := \{(x_i, y_i)\,|\, i \in [n]\}.
    \end{equation}
Here $D$ denotes the training dataset, and each $y_i \in \{-1,1\}$ is the ground-truth label prepared for the classification task.

Oppositely, we also wish to evaluate the model's ability to perform well on unseen or more general tasks.  
This motivates the definition of a new concept:

\begin{definition}[Out-of-sample Error]
The \emph{out-of-sample error} $E_{\rm out}$ measures how well a model generalizes beyond the training data.
\end{definition}

In the previous example, $E_{\rm out}$ can be expressed as
\begin{equation}
    E_{\rm out} = \mathbb{P}\big(h(x) \neq y\big)
    = \mathbb{E}_{(x,y)\sim p_{xy}}\!\left[\mathbb{1}\big(h(x) \neq y\big)\right],
\end{equation}
where $(x,y) \sim p_{xy}$ are i.i.d. samples drawn from the true underlying data distribution.

To quantitatively capture the generalization ability of a model, we define the following notion:

\begin{definition}[Generalization Error]
The \emph{generalization error} is defined as the difference between the out-of-sample and in-sample errors:
\begin{equation}
    E_{\rm out}(h) - E_{\rm in}(h).
\end{equation}
\end{definition}

\begin{remark}
In practice, we often introduce regularization terms to prevent $E_{\rm in}$ from becoming excessively small, which can help improve generalization performance.
\end{remark}

Our ultimate goal is to bound the quantity $E_{\rm out}(h) - E_{\rm in}(h)$.  
However, it remains an open question whether such a bound can be feasibly obtained for arbitrary hypotheses.
\begin{figure}[H]
    \centering
    \includegraphics[scale=0.8]{../../tikz/6/1.pdf}
    \caption{An example where $E_{\rm in}=0$ but $E_{\rm out}$ remains large.}
\end{figure}

It is crucial to acknowledge that, in essence, no finite dataset can guarantee a definitive prediction of future outcomes.  
The optimization of a model solely based on minimizing $E_{\rm in}(h)$ does not ensure a small generalization gap $E_{\rm out}(h) - E_{\rm in}(h)$, which may in fact be arbitrarily large.  
\textbf{However}, according to probabilistic learning theory, we can establish that with high probability $(1 - \delta)$, where $\delta \ll 1$, the generalization error is bounded by a small quantity $\varepsilon(\mathcal{H}, n, \delta)$ depending on the hypothesis space $\mathcal{H}$, the number of samples $n$, and the confidence parameter $\delta$:  
\begin{equation}
    \mathbb{P}\!\left(
        |E_{\rm out}(h) - E_{\rm in}(h)| \le \varepsilon(\mathcal{H}, n, \delta)
    \right) \ge 1 - \delta.
\end{equation}
This forms the cornerstone of statistical learning theory: while perfect certainty is unattainable, we can quantify uncertainty and guarantee performance with high probability.

The theorem above leads to the framework known as the \textbf{Probably Approximately Correct (PAC) Learning Theory}.  
Before introducing the PAC formulation in detail, we first review several fundamental mathematical tools.

\begin{theorem}[Hoeffding's Inequality]
Let $X_1, \dots, X_n$ be independent random variables such that each $X_i$ is almost surely bounded, $X_i \in [a_i, b_i]$.  
Define the empirical mean $\bar{X} = \frac{1}{n}\sum_{i=1}^n X_i$. Then, for any $\varepsilon > 0$, we have
\begin{equation}
    \mathbb{P}\!\left(|\bar{X} - \mathbb{E}\bar{X}| \ge \varepsilon\right)
    \le 2\exp\left(-\frac{2n^2\varepsilon^2}{\sum_{i=1}^n (b_i - a_i)^2}\right).
\end{equation}
\end{theorem}
\begin{proof}
    Let $S = \sum_{i=1}^n X_i$. We first bound the upper tail $\mathbb{P}(S - \mathbb{E}S \ge t)$; the lower tail is analogous and the two-sided bound follows by union.
    
    By Markov's inequality, for any $\lambda > 0$,
    \begin{equation}
    \mathbb{P}(S - \mathbb{E}S \ge t)
    = \mathbb{P}\!\left(e^{\lambda(S - \mathbb{E}S)} \ge e^{\lambda t}\right)
    \le e^{-\lambda t}\,\mathbb{E}\,e^{\lambda(S - \mathbb{E}S)}.
    \end{equation}
    
    Because $X_1,\dots,X_n$ are independent,
    \begin{equation}
    \mathbb{E}\,e^{\lambda(S - \mathbb{E}S)}
    = \prod_{i=1}^n \mathbb{E}\,e^{\lambda(X_i - \mathbb{E}X_i)}.
    \end{equation}
    
    For each bounded random variable $Y \in [a,b]$ with $\mathbb{E}Y = \mu$, Hoeffding's lemma gives
    \begin{equation}
    \mathbb{E}\,e^{\lambda(Y - \mu)} \le \exp\left(\frac{\lambda^2(b - a)^2}{8}\right).
    \end{equation}
    
    Applying this to each $X_i$ yields
    \begin{equation}
    \mathbb{E}\,e^{\lambda(S - \mathbb{E}S)}
    \le \exp\left(\frac{\lambda^2}{8}\sum_{i=1}^n (b_i - a_i)^2\right).
    \end{equation}
    
    Hence,
    \begin{equation}
    \mathbb{P}(S - \mathbb{E}S \ge t)
    \le \exp\left(-\lambda t + \frac{\lambda^2}{8}\sum_{i=1}^n (b_i - a_i)^2\right).
    \end{equation}
    
    Optimizing over $\lambda > 0$, we find
    \begin{equation}
    \lambda^\star = \frac{4t}{\sum_{i=1}^n (b_i - a_i)^2},
    \end{equation}
    which gives
    \begin{equation}
    \mathbb{P}(S - \mathbb{E}S \ge t)
    \le \exp\left(-\frac{2t^2}{\sum_{i=1}^n (b_i - a_i)^2}\right).
    \end{equation}
    
    By symmetry,
    \begin{equation}
    \mathbb{P}(S - \mathbb{E}S \le -t)
    \le \exp\left(-\frac{2t^2}{\sum_{i=1}^n (b_i - a_i)^2}\right).
    \end{equation}
    
    Combining both tails gives
    \begin{equation}
    \mathbb{P}\!\left(|S - \mathbb{E}S| \ge t\right)
    \le 2\exp\left(-\frac{2t^2}{\sum_{i=1}^n (b_i - a_i)^2}\right).
    \end{equation}
    
    Finally, setting $t = n\varepsilon$ and noting $\bar X = S/n$, we obtain
    \begin{equation}
    \mathbb{P}\!\left(|\bar X - \mathbb{E}\bar X| \ge \varepsilon\right)
    \le 2\exp\left(-\frac{2n^2\varepsilon^2}{\sum_{i=1}^n (b_i - a_i)^2}\right).
    \end{equation}
\end{proof}
Now, for a \textbf{given and fixed} hypothesis $h$, we obtain the following result:

\begin{theorem}[First Bound of PAC Learning]
For any $\varepsilon > 0$,
\begin{equation}
    \mathbb{P}\!\left(E_{\rm out}(h) - E_{\rm in}(h) \ge \varepsilon\right)
    \le \exp(-2n\varepsilon^2).
\end{equation}
\end{theorem}

Here, $E_{\rm in}(h)$ serves as the empirical mean of the random variables  
$\mathbb{1}\!\big(h(x_i) \neq y_i\big)$, as defined previously.  
In other words, each indicator $\mathbb{1}\!\big(h(x_i) \neq y_i\big)$ represents a bounded random variable in $[0,1]$,  
and Hoeffding’s inequality ensures that the deviation between the empirical error $E_{\rm in}(h)$  
and the expected error $E_{\rm out}(h)$ decreases exponentially with the number of samples $n$.
\begin{proof}
    By taking expectation over the definition of $E_{\rm in}$, we have
    \begin{align*}
        \mathbb{E}[\bar{X}]
        &= \frac{1}{n}\sum_{i=1}^n \mathbb{E}[X_i] \\
        &= \frac{1}{n}\sum_{i=1}^n \mathbb{E}_{(x_i,y_i)\sim p_{xy}}
           \big[\mathbb{1}(h(x_i)\neq y_i)\big] \\
        &= \frac{1}{n}\sum_{i=1}^n E_{\rm out}(h)
        = E_{\rm out}(h).
    \end{align*}
    The third equality follows directly from the definition of $E_{\rm out}(h)$.  
    Applying Hoeffding’s inequality, note that each random variable $\mathbb{1}(h(x_i)\neq y_i)$ is bounded in $\{0,1\}$,  
    so $(b_i - a_i)^2 = 1$ for all $i$, and the denominator $\sum_{i=1}^n (b_i - a_i)^2 = n$ cancels accordingly.  
    This completes the proof.
    \end{proof}
    
    However, there is a fundamental limitation in the condition of this first PAC bound:  
    the hypothesis $h$ must be fixed \emph{before} observing the training data.  
    In practice, the learning algorithm selects $h$ based on the dataset $D$,  
    which introduces statistical dependence between $h$ and $D$. (Parameters of model $h$ is determined by the chosen training dataset $D$) 
    This dependence violates the independence assumption required by Hoeffding’s inequality,  
    and therefore the bound no longer holds directly once $h$ is data-dependent.

    To overcome this limitation, we turn to analyzing the \textbf{union bound} over the hypothesis space $\mathcal{H}$,  
    which removes the dependency on a single fixed hypothesis chosen after observing the training data.  
    We first assume that $\mathcal{H}$ is a finite set, which actually is not always satisfied in practice  
    (e.g., the space of linear models has infinite cardinality, $|\mathrm{GL}_n| = \infty$).  
    That is, let $\mathcal{H} = \{h_1, \dots, h_M\}$.  
    
    In the worst case, we have
    \begin{align*}
        \mathbb{P}\!\left(\exists\, h \in \mathcal{H} : E_{\rm out}(h) - E_{\rm in}(h) \ge \varepsilon \right)
        &= \mathbb{P}\!\left(\bigcup_{i=1}^M \{E_{\rm out}(h_i) - E_{\rm in}(h_i) \ge \varepsilon\}\right) \\
        &\!\!\!\!\!\!\!\!\!\underbrace{\le}_{\text{Union Bound}} \sum_{i=1}^M
           \mathbb{P}\!\left(E_{\rm out}(h_i) - E_{\rm in}(h_i) \ge \varepsilon\right) \\
        &\le M \cdot \exp(-2n\varepsilon^2).
    \end{align*}
    
    Hence, we obtain the following result:
    
    \begin{theorem}[First Practical Bound of PAC Learning]
    With probability at least $1-\delta$,  
    \begin{equation}
        E_{\rm out}(h_i) - E_{\rm in}(h_i)
        < \sqrt{\frac{1}{2n}\log\frac{M}{\delta}}.
    \end{equation}
    \label{theo-6.3}
    \end{theorem}
\begin{remark}
    \begin{equation}
        \delta:=M\cdot\exp(-2n\varepsilon^2)\quad\Leftrightarrow \quad \varepsilon=\sqrt{\frac{1}{2n}\log\frac{M}{\delta}}
    \end{equation}
\end{remark}
    This result reveals a fundamental trade-off between model complexity and generalization.  
    As the hypothesis space $\mathcal{H}$ becomes richer, the total number of possible hypotheses $M$ increases.  
    A larger $M$ expands the logarithmic term $\log\frac{M}{\delta}$, thereby loosening the bound on $E_{\rm out}-E_{\rm in}$.  
    In other words, even though a highly expressive model may achieve $E_{\rm in}=0$ by perfectly fitting the training data,  
    its generalization gap can still remain large.  
    To maintain reliable generalization for such expressive models, one must increase the sample size $n$ accordingly.  
    This highlights the essential balance between data volume and model capacity in learning theory.

    Now consider an infinite hypothesis space $\mathcal{H}$, for example the reproducing kernel Hilbert space (RKHS)  
    of linear classifiers $\{h : h(x) = w^\top x + b\}$.  
    We claim that the bound provided in the theorem above becomes an extreme overestimation in such cases.  
    Intuitively, the union bound treats all hypotheses as completely distinct,  
    even if many of them make identical predictions on the dataset.  
    For example, the following three classifiers $h_1$, $h_2$, and $h_3$ may produce exactly the same classification results,  
    yet the union bound still counts them as three separate events,
    thereby inflating the estimated probability by a factor of three.  
\begin{figure}[H]
\centering
\includegraphics[scale=0.9]{../../tikz/6/2.pdf}

\end{figure}
\begin{note}
    A natural idea is to quotient out an equivalence relation on $\mathcal{H}$ so that its effective measure decreases.
This leads us to the concept of \emph{dichotomy}, where we seek an isomorphism that maps hypotheses producing identical classification results to the same image, thereby identifying them within a common kernel space.
By quotienting out the equivalence classes induced by $\ker$, which acts as a normal subgroup,
we obtain a new structure—still a group (or more appropriately here, a reduced hypothesis space)—
a space of smaller measure, serving as a kind of \emph{renormalized} hypothesis space.
\end{note}

Assume a binary classification setting
\begin{equation}
    y \in \{-1,+1\}, \qquad x \in \mathcal{X},
\end{equation}
with hypothesis space
\begin{equation}
    \mathcal{H} = \{h:\mathcal{X}\to\{-1,+1\}\}.
\end{equation}

Given $n$ samples
\begin{equation}
    x_1,x_2,\dots,x_n \in \mathcal{X},
\end{equation}
for any $h\in\mathcal{H}$ we obtain an $n$–tuple of labels
\begin{equation}
    \bigl(h(x_1),h(x_2),\dots,h(x_n)\bigr) \in \{-1,+1\}^n,
\end{equation}
which is called a \emph{dichotomy} of the set $\{x_1,\dots,x_n\}$ induced by $\mathcal{H}$.

The growth function $m_{\mathcal{H}}(n)$ is defined as:
\begin{definition}[Growth Function]
    \begin{equation}
        m_{\mathcal{H}}(n)
        = \max_{x_1,\dots,x_n \in \mathcal{X}}
          \Bigl|\bigl\{(h(x_1),\dots,h(x_n)) : h\in\mathcal{H}\bigr\}\Bigr|,
    \end{equation}
\end{definition}
i.e., the maximum number of distinct dichotomies that $\mathcal{H}$ can implement
on any set of $n$ points. It measures the effective size of the hypothesis space
$\mathcal{H}$.
\begin{note}
    The growth function can be interpreted as a quotient construction.  
    Define the evaluation map
    \begin{equation}
        \Phi_n : \mathcal{H} \to \{-1,+1\}^n, \qquad
        \Phi_n(h) = (h(x_1), h(x_2), \dots, h(x_n)).
    \end{equation}
    This map induces an equivalence relation on $\mathcal{H}$:
    \begin{equation}
        h_1 \sim h_2 \quad \Leftrightarrow \quad h_1(x_i) = h_2(x_i), \ \forall i = 1, \dots, n.
    \end{equation}
    Each equivalence class corresponds to one unique dichotomy,  
    so two hypotheses are considered identical if they generate the same labeling  
    on the given sample set. Therefore, the growth function can be written as
    \begin{equation}
        m_{\mathcal{H}}(n) = \bigl|\mathcal{H} / \ker(\Phi_n)\bigr|,
    \end{equation}
    representing the number of equivalence classes of $\mathcal{H}$ under this relation.  
    Intuitively, this quotient construction mods out by the kernel of the evaluation map,  
    focusing only on the empirically distinguishable behaviors of hypotheses on $n$ samples.
    \end{note}
    Since there are at most $2^n$ different $\pm1$ labelings on $n$ points, we always have
\begin{equation}
    m_{\mathcal{H}}(n) \le 2^n.
\end{equation}

Now let $S = \{x_1,\dots,x_n\} \subset \mathcal{X}$ be a finite subset.
We say that $\mathcal{H}$ \emph{shatters} $S$ if it can realize all possible
$\pm1$ assignments on $S$, that is,
\begin{equation}
    \mathcal{H}(x_1,\dots,x_n) = \{-1,+1\}^n.
\end{equation}
Equivalently, for every labeling $(y_1,\dots,y_n) \in \{-1,+1\}^n$
there exists some $h \in \mathcal{H}$ such that $h(x_i) = y_i$ for all $i=1,\dots,n$.
\begin{example}
    Consider three sample points $x_1, x_2, x_3$ and a hypothesis set $\mathcal{H}$
    that induces the following dichotomies:
    \begin{equation}
        \mathcal{H}(x_1,x_2,x_3)
        = \{ (+1,-1,-1),\; (-1,+1,-1),\; (-1,+1,+1) \}.
    \end{equation}
Determine which subsets of $\{x_1, x_2, x_3\}$ can be \emph{shattered} by $\mathcal{H}$.
\end{example}
\begin{solution}
    \begin{itemize}
        \item For any single point $\{x_i\}$, both $+1$ and $-1$ appear in the projections
        of $\mathcal{H}(x_1,x_2,x_3)$, so every singleton subset is shattered.
    
        \item For pairs of points:
        \begin{align*}
            \{x_1, x_2\} &\Rightarrow \{(+1,-1),\, (-1,+1)\} \quad
            \text{(missing $(+1,+1)$ and $(-1,-1)$) $\Rightarrow$ not shattered;} \\
            \{x_1, x_3\} &\Rightarrow \{(+1,-1),\, (-1,+1),\, (-1,-1)\} \quad
            \text{(missing $(+1,+1)$) $\Rightarrow$ not shattered;} \\
            \{x_2, x_3\} &\Rightarrow \{(-1,-1),\, (+1,-1),\, (+1,+1)\} \quad
            \text{(all $2^2=4$ combinations not present) $\Rightarrow$ not shattered.}
        \end{align*}
    
        \item For the full set $\{x_1,x_2,x_3\}$, since only $3<2^3=8$ dichotomies are realized,
        it is clearly not shattered.
    \end{itemize}
    
    Hence, $\mathcal{H}$ can only shatter the \emph{singleton subsets}
    $\{x_1\}$, $\{x_2\}$, and $\{x_3\}$, but no pair or triple of points.
\end{solution}
\begin{example}
    Let $\mathcal{X} = \mathbb{R}^2$, and consider the hypothesis set of all linear classifiers
    \begin{equation}
        \mathcal{H} = \{\, h_{w,b}(x) = \operatorname{sign}(w^\top x + b) \mid w \in \mathbb{R}^2,\, b \in \mathbb{R} \,\}.
    \end{equation}
    Determine the value of $m_{\mathcal{H}}(3)$, the maximum number of dichotomies that $\mathcal{H}$
    can produce on any $3$ points in $\mathbb{R}^2$.
\end{example}
\begin{solution}
    

A separating hyperplane in $\mathbb{R}^2$ is a line of the form $w^\top x + b = 0$.  
The sign of $w^\top x + b$ assigns each point to one of two half-spaces.
 
For any given labeling $(y_1,y_2,y_3)\in\{-1,+1\}^3$,  
we can always find some line that places the $+1$ points on one side and the $-1$ points on the other,  
as long as the three points are not collinear.  
This is because two points determine a line, and the third point’s label can always be satisfied
by shifting or rotating that line slightly.

Hence every possible assignment of $\pm1$ labels to the three points is realizable:
\begin{equation}
    \mathcal{H}(x_1,x_2,x_3)
    = \{-1,+1\}^3.
\end{equation}

Therefore, the growth function value is
\begin{equation}
    m_{\mathcal{H}}(3) = 2^3 = 8.
\end{equation}
\end{solution}
\begin{remark}
    More examples have been presented in class, but I'm not goint to note these in our lecture.
\end{remark}
$m_{\mathcal H}(n)$ quantified the ability of a model.
\section{VC-Dim}

\begin{definition}[Vapnik--Chervonenkis Dimension]
The \emph{VC dimension} $\mathrm{dvc}(\mathcal{H})$ is the largest integer $n$ such that
\begin{equation}
    m_{\mathcal{H}}(n) = 2^n.
\end{equation}
\end{definition}

\begin{remark}
Vapnik is the inventor of the Support Vector Machine (SVM) and one of the founders of
statistical learning theory.
\end{remark}

The VC dimension measures the \emph{capacity} or \emph{expressive power} of a hypothesis class.
Formally, it is the largest integer $n$ such that $\mathcal{H}$ can realize all
$2^n$ possible dichotomies on some set of $n$ points.
Equivalently, it quantifies how complex the decision boundaries in $\mathcal{H}$ can be:
a larger VC dimension means $\mathcal{H}$ can represent more distinct labelings of data,
i.e.\ greater capacity.

\smallskip
It turns out that if $\mathrm{dvc}(\mathcal{H})$ is finite,
then $m_{\mathcal{H}}(n)$ can be bounded in terms of $\mathrm{dvc}(\mathcal{H})$.

\begin{theorem}[(Sauer--)Shelah's Lemma]
Let $\mathcal{H}$ be a class of binary-valued functions on a domain $\mathcal{X}$.
Define its growth function as
\begin{equation}
    m_{\mathcal{H}}(n)
    = \max_{\substack{S \subseteq \mathcal{X} \\ |S| = n}}
      \bigl|\{\, (h(x))_{x \in S} : h \in \mathcal{H} \}\bigr|.
\end{equation}
If $\mathrm{dvc}(\mathcal{H}) = d < \infty$, then for all $n \ge 0$,
\begin{equation}
    m_{\mathcal{H}}(n)
    \le \sum_{i=0}^{d} \binom{n}{i}\sim O(n^{d})
\end{equation}
\end{theorem}
\begin{proof}
    We proceed by induction on $n$. The base case is trivial and we proved that in the next proof for completeness.
    
    \textbf{Step 1: Divide the hypothesis class.}
    Fix $x_1, \dots, x_n \in \mathcal{X}$, and define
    \begin{align}
        \mathcal{H}_1(x_1,\dots,x_{n-1})
        &= \{ (h(x_1), \dots, h(x_{n-1})) \mid h \in \mathcal{H},\, h(x_n)=+1 \}, \\
        \mathcal{H}_2(x_1,\dots,x_{n-1})
        &= \{ (h(x_1), \dots, h(x_{n-1})) \mid h \in \mathcal{H},\, h(x_n)=-1 \}.
    \end{align}
    Then
    \begin{equation}
        \mathcal{H}(x_1,\dots,x_n)
        = \{ (\mathbf{v}, +1) \mid \mathbf{v} \in \mathcal{H}_1(x_1,\dots,x_{n-1}) \}
          \cup
          \{ (\mathbf{v}, -1) \mid \mathbf{v} \in \mathcal{H}_2(x_1,\dots,x_{n-1}) \}.
    \end{equation}
    
    \textbf{Step 2: Counting distinct dichotomies.}
    If $\mathbf{v}$ belongs to both $\mathcal{H}_1$ and $\mathcal{H}_2$, then $\mathcal{H}$ realizes both labelings $(\mathbf{v}, +1)$ and $(\mathbf{v}, -1)$.  
    Hence,
    \begin{equation}
        |\mathcal{H}(x_1,\dots,x_n)|
        \le |\mathcal{H}_1(x_1,\dots,x_{n-1})|
           + |\mathcal{H}_2(x_1,\dots,x_{n-1})|.
    \end{equation}
    
    \textbf{Step 3: Bounding the VC dimension.}
    If $\mathrm{dvc}(\mathcal{H}) = d$, then $\mathrm{dvc}(\mathcal{H}_1) \le d$ and $\mathrm{dvc}(\mathcal{H}_2) \le d-1$.  
    Otherwise, if $\mathcal{H}_2$ could shatter a subset of size $d$, then $\mathcal{H}$ could shatter that subset together with $x_n$, contradicting $\mathrm{dvc}(\mathcal{H}) = d$.
    
    \textbf{Step 4: Apply the induction hypothesis.}
    By the inductive assumption,
    \begin{equation}
        |\mathcal{H}_1(x_1,\dots,x_{n-1})| \le \sum_{i=0}^{d} \binom{n-1}{i},
        \qquad
        |\mathcal{H}_2(x_1,\dots,x_{n-1})| \le \sum_{i=0}^{d-1} \binom{n-1}{i}.
    \end{equation}
    Therefore,
    \begin{align}
        |\mathcal{H}(x_1,\dots,x_n)|
        &\le \sum_{i=0}^{d} \binom{n-1}{i}
           + \sum_{i=0}^{d-1} \binom{n-1}{i} \\
        &= \sum_{i=0}^{d} \left( \binom{n-1}{i} + \binom{n-1}{i-1} \right)
           = \sum_{i=0}^{d} \binom{n}{i},
    \end{align}
    where the last equality uses Pascal's identity.  
    Hence the claim holds for all $n$.
    \end{proof}
    \begin{tcolorbox}
        \vspace{0.5em}
            \color{red!50!black}$\circledast   $\quad\emph{{The following proof is included out of personal interest and is not part of the course material. 
            Readers who wish to focus solely on the course content may safely skip to the next part.}}
        \vspace{0.5em}
    \end{tcolorbox}
\begin{proof}
    Here we provide another proof.
    
    Fix $n$ and choose $S \subseteq \mathcal{X}$ with $|S|=n$
    such that $m_{\mathcal{H}}(n)=|\mathcal{H}|_S|$, where
    \begin{equation}
        \mathcal{H}|_S
        := \{\, (h(x))_{x\in S} : h\in\mathcal{H} \}
    \end{equation}
    is the set of labelings realized on $S$.
    We prove by induction on $n$ that
    \begin{equation}
        |\mathcal{H}|_S|
        \le \sum_{i=0}^{d} \binom{n}{i}.
    \end{equation}
    
    \textbf{Base cases.}
    If $n=0$, then $|\mathcal{H}|_S|=1$ and
    \begin{equation}
        \sum_{i=0}^{d} \binom{0}{i} = 1,
    \end{equation}
    so the claim holds.
    If $d=0$, then no single point can be labeled in two different ways,
    hence $|\mathcal{H}|_S|=1$ for all $n$, while
    \begin{equation}
        \sum_{i=0}^{0} \binom{n}{i} = 1,
    \end{equation}
    so the bound again holds.
    
    \textbf{Inductive step.}
    Assume the statement true for all pairs $(n',d')$ with $n' < n$
    and the same $d$. Pick some $x \in S$ and set
    \begin{equation}
        S' := S \setminus \{x\}, \qquad |S'| = n-1.
    \end{equation}
    For each $h\in\mathcal{H}$, let $h|_{S'}$ be its restriction to $S'$.
    Define
    \begin{equation}
        \mathcal{A}
        := \{\, h|_{S'} : h\in\mathcal{H} \},
    \end{equation}
    \begin{equation}
        \mathcal{B}
        := \bigl\{\, h|_{S'} : h\in\mathcal{H},\;
            \exists\,h' \in \mathcal{H}
            \text{ with } h'(x)\neq h(x),\ h'|_{S'}=h|_{S'} \bigr\}.
    \end{equation}
    Thus, $\mathcal{A}$ is the set of all restrictions on $S'$,
    and $\mathcal{B}$ collects those restrictions for which
    $\mathcal{H}$ contains at least two hypotheses that agree on $S'$
    but disagree on $x$.
    
    Every labeling of $S$ comes from choosing a restriction in $\mathcal{A}$,
    and then choosing the value at $x$.
    Labelings for which the value at $x$ is \emph{forced} (no alternative
    hypothesis with the same restriction) contribute at most $|\mathcal{A}|$
    patterns; labelings for which the value at $x$ can be flipped
    (i.e.\ restrictions in $\mathcal{B}$) contribute at most $|\mathcal{B}|$
    additional patterns.
    Hence
    \begin{equation}
        |\mathcal{H}|_S|
        \le |\mathcal{A}| + |\mathcal{B}|.
    \end{equation}
    
    Note that $|\mathcal{A}| \le m_{\mathcal{H}}(n-1)$, because $\mathcal{A}$
    is a set of labelings on $S'$ of size $n-1$.
    To bound $|\mathcal{B}|$, define a new hypothesis class
    \begin{equation}
        \mathcal{H}_x
        := \bigl\{\, h|_{S'} : h\in\mathcal{H},\;
            \exists\,h' \in \mathcal{H}
            \text{ with } h'(x)\neq h(x),\ h'|_{S'}=h|_{S'} \bigr\}.
    \end{equation}
    Then $\mathcal{B}$ is precisely $\mathcal{H}_x$ viewed as a set of labelings on $S'$,
    so $|\mathcal{B}| \le m_{\mathcal{H}_x}(n-1)$.
    
    We claim that
    \begin{equation}
        \mathrm{dvc}(\mathcal{H}_x) \le d-1.
    \end{equation}
    Otherwise, suppose $\mathcal{H}_x$ shatters some $T \subseteq S'$ with $|T|=d$.
    By the definition of $\mathcal{H}_x$, for every labeling $y_T$ on $T$
    and either choice of label $y_x \in \{-1,+1\}$ at $x$, there exists $h\in\mathcal{H}$
    such that
    \begin{equation}
        h|_{T} = y_T, \qquad h(x) = y_x.
    \end{equation}
    Hence $\mathcal{H}$ can realize all labelings on $T \cup \{x\}$,
    so $\mathcal{H}$ shatters a set of size $d+1$, contradicting
    $\mathrm{dvc}(\mathcal{H}) = d$.
    Thus the claim holds.
    
    By the induction hypothesis applied to $\mathcal{H}$ on $n-1$ points
    and to $\mathcal{H}_x$ whose VC dimension is at most $d-1$, we obtain
    \begin{equation}
        |\mathcal{A}| \le \sum_{i=0}^{d} \binom{n-1}{i},
        \qquad
        |\mathcal{B}| \le \sum_{i=0}^{d-1} \binom{n-1}{i}.
    \end{equation}
    Therefore,
    \begin{align}
        |\mathcal{H}|_S|
        &\le \sum_{i=0}^{d} \binom{n-1}{i}
           + \sum_{i=0}^{d-1} \binom{n-1}{i} \\
        &= \sum_{i=0}^{d} \left( \binom{n-1}{i} + \binom{n-1}{i-1} \right) \\
        &= \sum_{i=0}^{d} \binom{n}{i}.
    \end{align}
    \end{proof}
    \begin{note}
        The inductive structure of the proof for the Sauer--Shelah Lemma is deeply reminiscent
        of strategies used in \emph{mathematical logic}, particularly in proving
        \emph{completeness theorems} that connect syntactic derivability and semantic truth.
        In both settings, we start by imposing an ordering (or enumeration) on a collection
        of objects—whether they are formulas, sentences, or finite subsets of $\mathcal{X}$—so that
        induction can be applied coherently.
        
        Just as in logic we construct models step by step according to a fixed sequence
        of formulas (often by extending consistent sets in a well-ordered fashion),
        here we inductively analyze the growth of $\mathcal{H}$ by removing or reintroducing one
        element of the sample set at a time.  
        The key similarity is structural: both arguments rely on well-founded recursion over
        a totally ordered domain to control how local extensions (adding one more formula or one more
        point) affect global consistency or combinatorial growth.
        \end{note}
        
    \begin{remark}
    Plugging this bound into the standard uniform convergence inequality
    for a finite hypothesis class, we obtain a generalization bound of the form
    \begin{equation}
        \varepsilon
        \;=\;
        \sqrt{\frac{1}{2n}\,\log \frac{M}{\delta}}\sim O\left(\frac{\mathrm{dvc}\log n}{n}\right)^{1/2}
    \end{equation}
    Thus, when $\mathrm{dvc}(\mathcal{H})$ is finite, the estimation error
    $\varepsilon$ decreases as $n$ grows, with the complexity of the class
    entering only through the polynomial factor controlled by $d$.
    \end{remark}
    It can be rigorously shown that:

    \begin{theorem}[VC Generalization Bound]
    With probability at least $1-\delta$, the following inequality holds:
    \begin{equation}
        E_{\rm out}(h_i) - E_{\rm in}(h_i)
        < \sqrt{\frac{8}{n}\log\frac{4m_{\mathcal{H}}(2n)}{\delta}}.
    \end{equation}
    Moreover,
    \begin{equation}
        E_{\rm out}(h_i) - E_{\rm in}(h_i)
        = O\!\left(\sqrt{\frac{\mathrm{d_{VC}}\log n}{n}}\right)
        \xrightarrow[n\to\infty]{} 0.
    \end{equation}
    \end{theorem}
    
    Although the VC generalization bound is relatively loose in practice,  
    its importance lies in converting an inherently infinite hypothesis space into a finitely bounded one.  
    This theorem establishes a concrete analytical foundation for the entire field of machine learning,  
    demonstrating that with sufficient data, generalization is not merely empirical, but theoretically guaranteed.
\end{document}